\section*{Chapter \ref{ch:g6:water-underground} --- Water Underground}

\begin{solution}{exo:g6:water-underground:1}
\omterm{def:g6:water-underground:permeability}{Permeable}: \omterm{def:g3:stones-pebbles-sand:sand}{sand}, \omterm{def:g3:stones-pebbles-sand:pebble}{gravel}, chalk. \omterm{def:g6:water-underground:permeability}{Impermeable}: clay, granite without cracks,
shale.
\end{solution}

\begin{solution}{exo:g6:water-underground:2}
(a) \omterm{def:g6:water-underground:water-table}{water table}; (b) \omterm{def:g6:water-underground:aquifer}{aquifer}; (c) \omterm{def:g6:water-underground:spring}{spring}; (d) \omterm{def:g6:water-underground:spring}{well}.
\end{solution}

\begin{solution}{exo:g6:water-underground:3}
(a) It evaporates back into the air; it runs off over the surface into
streams; it soaks into the ground (\omterm{def:g6:water-underground:infiltration}{infiltration}). (b) On the car park: its
hard surface lets almost no water soak in, so the water runs off. In the
meadow, the soil and the grass let much of it soak in.
\end{solution}

\begin{solution}{exo:g6:water-underground:4}
(a) False: clay is very \omterm{def:g6:water-underground:porosity}{porous} but \omterm{def:g6:water-underground:permeability}{impermeable}, because its \omterm{def:g6:water-underground:porosity}{pores} are tiny.
(b) True. (c) False: it follows the shape of the land, more gently, and it
goes up and down with the seasons. (d) True: about 30~\% of all \omterm{def:g2:water-on-earth:salt-fresh}{fresh water}
is \omterm{def:g6:water-underground:groundwater}{groundwater}, against about a quarter of one percent in the \omterm{def:g1:landscapes-around-us:water}{lakes}.
\end{solution}

\begin{solution}{exo:g6:water-underground:5}
$10 \times 25 \div 100 = 2.5$: you can pour 2.5~litres before it overflows.
$10 \times 22 \div 100 = 2.2$: 2.2~litres would drain out through the hole.
\end{solution}

\begin{solution}{exo:g6:water-underground:6}
(a) $2 \times 1 \times 1 = 2$~cubic metres, that is 2000~litres.
(b) $2000 \times 20 \div 100 = 400$~litres. (c) $2000 \times 18 \div 100 =
360$~litres drain out; $400 - 360 = 40$~litres stay stuck.
\end{solution}

\begin{solution}{exo:g6:water-underground:7}
(a) The \omterm{def:g6:water-underground:spring}{spring} is at about 203~m, the \omterm{def:g1:landscapes-around-us:water}{river} at about 132~m. (b) The \omterm{def:g6:water-underground:spring}{spring}
is about 70~m above the \omterm{def:g1:landscapes-around-us:water}{river}: its water runs down the slope into the \omterm{def:g1:landscapes-around-us:water}{river}.
(c) Inside the chalk hillside the \omterm{def:g6:water-underground:water-table}{water table} is below the ground; it meets
the surface only where the base of the chalk, resting on the \omterm{def:g6:water-underground:permeability}{impermeable}
shale, comes out on the slope.
\end{solution}

\begin{solution}{exo:g6:water-underground:8}
(a) After the rainy season: $401 - 336 = 65$~m; after the dry season:
$401 - 314 = 87$~m. (b) $336 - 314 = 22$~m.
\end{solution}

\begin{solution}{exo:g6:water-underground:9}
Deeper. The \omterm{def:g6:water-underground:water-table}{water table} follows the shape of the land, but more gently: at
the top of the \omterm{def:g1:landscapes-around-us:landform}{hill} it is higher than under \omterm{def:g6:water-underground:spring}{well}~A, but the ground rises
much more, so the water lies farther below the surface (on the section,
more than 110~m down).
\end{solution}

\begin{solution}{exo:g6:water-underground:10}
The \omterm{def:g6:water-underground:porosity}{pores} of clay are tiny: the water stays stuck to the flakes and cannot
flow through the clay towards the \omterm{def:g6:water-underground:spring}{well}. Clay is \omterm{def:g6:water-underground:porosity}{porous} but \omterm{def:g6:water-underground:permeability}{impermeable}, so
only a little water (20~litres out of 500) drains out.
\end{solution}

\begin{solution}{exo:g6:water-underground:11}
(a) $30.1 \div 0.26 \approx 116$: more than a hundred times more.
(b) $100 - 69 = 31$: about 31~\% is left, and \omterm{def:g6:water-underground:groundwater}{groundwater} alone is about 30
of them. (c) It is by far the largest store of liquid \omterm{def:g2:water-on-earth:salt-fresh}{fresh water}; the ice
is frozen and mostly far away, in the polar regions and on \omterm{def:g1:landscapes-around-us:landform}{mountains}. And
\omterm{def:g6:water-underground:groundwater}{groundwater} is found almost everywhere under our feet.
\end{solution}

\begin{solution}{exo:g6:water-underground:12}
Both \omterm{def:g6:water-underground:spring}{wells} draw on the same chalk \omterm{def:g6:water-underground:aquifer}{aquifer}. The powerful pump took water out
faster than the \omterm{def:g4:clouds-rain-snow:precipitation}{rain} refills the chalk, so the \omterm{def:g6:water-underground:water-table}{water table} went down all
around, below the bottom of the second \omterm{def:g6:water-underground:spring}{well}, which ran dry. The farmers
could agree to pump less, especially in summer, share one \omterm{def:g6:water-underground:spring}{well}, or deepen
the second \omterm{def:g6:water-underground:spring}{well}; above all they should not take more water than the \omterm{def:g4:clouds-rain-snow:precipitation}{rain}
brings back.
\end{solution}

\begin{solution}{pb:g6:water-underground:1}
\textbf{1.} The farm stands on chalk; under it lies shale. The chalk is
\omterm{def:g6:water-underground:permeability}{permeable}, the shale is not.

\textbf{2.} \omterm{def:g4:clouds-rain-snow:precipitation}{Rain} soaks into the chalk and stops on the shale. Where the
base of the chalk comes out on the slope, the \omterm{def:g6:water-underground:water-table}{water table} meets the ground
and the water comes out: that is the \omterm{def:g6:water-underground:spring}{spring}.

\textbf{3.} $401 - 336 = 65$~m below the ground.

\textbf{4.} The bottom of the \omterm{def:g6:water-underground:spring}{well} is at $401 - 100 = 301$~m, below the
\omterm{def:g6:water-underground:water-table}{water table} at 314~m: the \omterm{def:g6:water-underground:spring}{well} still reaches the water, by $314 - 301 =
13$~m.

\textbf{5.} $100 \times 50 \times 2 = 10\,000$~cubic metres.

\textbf{6.} $10\,000 \times 20 \div 100 = 2000$~cubic metres.

\textbf{7.} $10\,000 \times 18 \div 100 = 1800$~cubic metres drain out;
$2000 - 1800 = 200$~cubic metres stay stuck.

\textbf{8.} $1800 \times 1000 = 1\,800\,000$~litres.

\textbf{9.} $18 \times 1000 = 18\,000$~litres a day.

\textbf{10.} $18 \times 7 = 126$~cubic metres a week.

\textbf{11.} The \omterm{def:g4:clouds-rain-snow:precipitation}{rain} of the wet season soaks into the chalk and refills it
every year (as long as the farms do not pump more than the \omterm{def:g4:clouds-rain-snow:precipitation}{rain} brings back).

\textbf{12.} $1800 \div 18 = $~\textbf{100 days}: a little more than three
months without any \omterm{def:g4:clouds-rain-snow:precipitation}{rain}.
\end{solution}
